\subsection{Nomenclatura}

    En esta sección se establecen las convenciones de notación utilizadas a lo largo del manuscrito. Se sigue la convención de Einstein expuesta en \textit{Los Fundamentos de la Teoría de la Relatividad General} \cite{einstein1916fundamentos}.

    \subsubsection{Tensores contravariantes y covariantes}

        \begin{definicion}{Notación tensorial}{def:notacion_tensorial}
            En el formalismo tensorial de relatividad general, se adopta la siguiente \textit{notación tensorial}:
            \begin{itemize}
                \item $A^\mu$ denota un \textit{cuadrivector contravariante}.
                \item $A_\mu$ denota un \textit{cuadrivector covariante}.
                \item $A^{\mu\nu}$ denota un \textit{tensor de segundo orden contravariante}.
                \item $A_{\mu\nu}$ denota un \textit{tensor de segundo orden covariante}.
                \item ${A^{\mu}}_{\nu}$ denota un \textit{tensor de segundo orden mixto}.
                \item ${A^{\mu_1\mu_2...\mu_n}}_{\nu_1\nu_2...\nu_m}$ denota un \textit{tensor de orden $n+m$}.
            \end{itemize}
        \end{definicion}

        El orden de cada índice, ya sea contravariante o covariante es importante, por ello la separación de espacio entre superíndices y subíndices.

        \begin{observacion}{Índices griegos y latinos}{obs:indices_rango}
            Establecemos como convención que los índices griegos recorran los valores 0, 1, 2 y 3, y los índices latinos recorran los valores 1, 2 y 3. En este contexto, el índice 0 corresponde a la coordenada relacionada con el tiempo.
        \end{observacion}

    \subsubsection{Contracción de índices}

        \begin{definicion}{Contracción de índices}{def:contraccion_indices}
            Sean ${A^{\mu_1\mu_2...\mu_n}}_{\nu_1\nu_2...\nu_m}$ y ${B^{\nu_1\nu_2...\nu_m}}_{\rho_1\rho_2...\rho_n}$ dos tensores. La \textit{contracción de índices} entre ambos se define como:
            \begin{equation}
                {A^{\mu_1\mu_2...\mu_n}}_{\nu_1\nu_2...\nu_m} {B^{\nu_1\nu_2...\nu_m}}_{\rho_1\rho_2...\rho_n} = \sum_{\nu_1 = 0}^{3} \sum_{\nu_2 = 0}^{3} ... \sum_{\nu_m = 0}^{3} {A^{\mu_1\mu_2...\mu_n}}_{\nu_1\nu_2...\nu_m} {B^{\nu_1\nu_2...\nu_m}}_{\rho_1\rho_2...\rho_n} \label{eq:contraccion_indices}
            \end{equation}
            Los índices repetidos en ambos tensores se denominan \textit{índices mudos}, y los índices que no se repiten se denominan \textit{índices libres}.
        \end{definicion}

        Para evitar ambigüedades, se establecen las siguientes reglas:

        \begin{observacion}{Reglas de la convención de Einstein}{obs:reglas_einstein}
            \begin{itemize}
                \item Cada término de una expresión dada puede contener a lo sumo dos índices repetidos. Por ejemplo, el término $A^\mu B_\mu C^\mu$ es inválido, pero el término $A^\mu B_\mu C^\nu$ es válido.
                \item Los índices sólo se pueden repetir si son de distinto tipo. Por ejemplo, $A^\mu B_\mu$ es válido, pero $A^\mu B^\mu$ es inválido.
                \item Los índices libres deben serlo para cada término de la expresión. Por ejemplo, el término $A^\mu = {B^\mu}_{\nu}C^\nu + D^\mu E_\mu$ es inválido, pero el término $A^\mu = {B^\mu}_{\nu}C^\nu + D^\mu$ es válido.
            \end{itemize}
        \end{observacion}

    \subsubsection{Derivadas}

        \paragraph{Notación de derivadas espaciales}

            Se establecen las siguientes notaciones para derivadas parciales.

            \begin{definicion}{Derivadas parciales covariante y contravariante}{def:derivadas_parciales}
                Sea $x^\mu$ un sistema de coordenadas en el espacio-tiempo $\mathcal{M}$. La \textit{derivada parcial covariante} $\partial_\mu$ y la \textit{derivada parcial contravariante} $\partial^\mu$ se definen como:
                \begin{align}
                    \partial_\mu \equiv \frac{\partial}{\partial x^\mu} \label{eq:derivada_parcial_covariante} \\
                    \partial^\mu \equiv \frac{\partial}{\partial x_\mu} \label{eq:derivada_parcial_contravariante}
                \end{align}
            \end{definicion}

            Dado que se deriva con frecuencia, una notación más compacta es:

            \begin{definicion}{Notación de coma}{def:notacion_coma}
                Sea $A^\mu$ un cuadrivector contravariante definido en $\mathcal{M}$. La \textit{notación de coma} para derivadas parciales se define como:
                \begin{align}
                    A^{\mu, \nu} \equiv \partial^\nu A^\mu \label{eq:notacion_coma_contravariante} \\
                    {A^{\mu}}_{,\nu} \equiv \partial_\nu A^\mu \label{eq:notacion_coma_covariante}
                \end{align}
                En general, para $n$ derivadas parciales sucesivas:
                \begin{align}
                    A^{\nu, \mu_1, \mu_2, ..., \mu_n} \equiv \partial^{\mu_n} ... \partial^{\mu_2} \partial^{\mu_1} A^\nu \label{eq:notacion_coma_contravariante_n} \\
                    {A^{\nu}}_{,\mu_1, \mu_2, ..., \mu_n} \equiv \partial_{\mu_n} ... \partial_{\mu_2} \partial_{\mu_1} A^\nu \label{eq:notacion_coma_covariante_n}
                \end{align}
                Esta notación es válida para cualquier tensor.
            \end{definicion}

        \paragraph{Notación de derivadas temporales}

            En relatividad general, el tiempo coordenado es relativo.

